% 同济大学 计算机网络课程设计 实验报告
% 题目二：中小企业网络规划与设计 (V2)
% 基于 Cisco Packet Tracer Hub-and-Spoke 架构

\documentclass[12pt,a4paper]{ctexart}

% ============ 页面设置 ============
\usepackage[top=2.5cm, bottom=2.5cm, left=2.8cm, right=2.8cm]{geometry}
\usepackage{graphicx}
\usepackage{float}
\usepackage{hyperref}
\usepackage{amsmath}
\usepackage{amssymb}
\usepackage{enumitem}
\usepackage{listings}
\usepackage{xcolor}
\usepackage{booktabs}
\usepackage{longtable}
\usepackage{array}
\usepackage{caption}
\usepackage{subcaption}
\usepackage{fancyhdr}
\usepackage{titlesec}

% ============ 修复 fancyhdr 警告 ============
\setlength{\headheight}{14.5pt}

% ============ 代码块样式 ============
\definecolor{codebg}{rgb}{0.95,0.95,0.95}
\definecolor{codefg}{rgb}{0.1,0.1,0.3}
\definecolor{commentcolor}{rgb}{0.2,0.5,0.2}
\lstset{
    basicstyle=\small\ttfamily,
    backgroundcolor=\color{codebg},
    frame=single,
    rulecolor=\color{black!30},
    breaklines=true,
    postbreak=\mbox{\textcolor{red}{$\hookrightarrow$}\space},
    commentstyle=\color{commentcolor},
    keywordstyle=\color{blue},
    stringstyle=\color{orange},
    showstringspaces=false,
    numbers=left,
    numbersep=5pt,
    numberstyle=\tiny\color{gray},
    tabsize=2
}

% ============ 超链接设置 ============
\hypersetup{
    colorlinks=true,
    linkcolor=blue,
    urlcolor=cyan,
    citecolor=green
}

% ============ 页眉页脚 ============
\pagestyle{fancy}
\fancyhf{}
\fancyhead[L]{计算机网络课程设计}
\fancyhead[R]{题目二：中小企业网络规划与设计}
\fancyfoot[C]{\thepage}
\renewcommand{\headrulewidth}{0.4pt}

% ============ 标题格式 ============
\titleformat{\section}{\Large\bfseries}{}{0em}{}[]
\titleformat{\subsection}{\large\bfseries}{}{0em}{}[]
\titleformat{\subsubsection}{\normalsize\bfseries}{}{0em}{}[]

% ============ 图片路径 ============
\graphicspath{{Assets/v2/}}

\begin{document}

% ==================== 封面 ====================
\begin{titlepage}
\thispagestyle{empty}
\begin{center}
    \vspace*{1cm}

    {\Huge\textbf{同济大学}}\\[0.2cm]
    {\Huge\textbf{计算机科学与技术系}}\\[1.2cm]

    {\LARGE\textbf{计算机网络课程设计}}\\[0.3cm]
    {\LARGE\textbf{实验报告}}\\[1.2cm]

    
\includegraphics[width=0.4\textwidth]{../../Tongji.png}\\[1.2cm]

    \begin{large}
    \begin{tabular}{rl}
        \textbf{学\qquad 号：} & 2350998、2350988、2353598 \\[0.2cm]
        \textbf{姓\qquad 名：} & 庄子硕、何东峻、张云夏 \\[0.2cm]
        \textbf{专\qquad 业：} & 计算机科学与技术 \\[0.2cm]
        \textbf{授课老师：} & 田春岐 \\[0.2cm]
        \textbf{日\qquad 期：} & \today \\
    \end{tabular}
    \end{large}

    \vfill
\end{center}
\end{titlepage}

% ==================== 目录 ====================
\tableofcontents
\newpage

% ==================== 一、项目概述 ====================
\section{项目概述}

\subsection{项目背景}

随着企业信息化建设不断深入，中小企业对网络系统的依赖程度日益提高。企业内部办公自动化（OA）、邮件通信、文件共享、远程办公以及互联网访问等业务均需要稳定、安全、高效的网络环境作为支撑。

本项目针对某中小企业进行网络规划与设计。企业总部拥有约 1000 台 PC，并设有市场部、财务部、人事部、研发部、行政部、运营部、客服部共 7 个部门，同时在外省设有 3 个分公司。企业要求实现总部与分公司的互联互通，并能够安全地对外提供 WEB、邮件以及匿名 FTP 服务，同时保障内部 OA 系统的正常运行。

本方案基于 Cisco Packet Tracer 平台，采用\textbf{Hub-and-Spoke（中心辐射型）}网络拓扑架构，以中心路由器为核心，通过高速串行链路连接 ISP、分公司及总部网络。网络采用 \textbf{RIPv2 动态路由协议}实现路由自动学习，结合 \textbf{VLAN、ACL、VPN（IPSec）、NAT、DHCP} 等关键网络技术，构建安全、可靠、可扩展的企业网络环境。

\subsection{项目需求}

根据题目要求，网络系统需满足以下需求：

\begin{enumerate}[label=(\arabic*), leftmargin=2em]
    \item 公司总部拥有约 1000 台 PC；
    \item 企业共有 7 个部门，不同部门之间访问权限受限；
    \item 企业拥有 3 个跨省分公司；
    \item 企业内部部署 OA 系统；
    \item 企业可对外提供 WWW、匿名 FTP、邮件服务；
    \item FTP 服务仅允许内部员工访问；
    \item 所有终端均可访问互联网；
    \item 每个部门独立划分 VLAN；
    \item 总部与分公司之间通过 VPN 技术实现安全互联；
    \item 支持远程办公人员通过 VPN 接入企业内网；
    \item 网络需具备较高安全性、可靠性与扩展能力。
\end{enumerate}

% ==================== 二、可行性分析报告 ====================
\section{可行性分析报告}

\subsection{社会可行性分析}

企业信息化网络建设能够显著提升企业办公效率，促进企业内部信息共享与资源协同。通过构建统一的企业网络平台，可以实现跨部门协作、远程办公、数据共享以及业务系统集中管理。

VPN 技术能够有效保障总部与分公司之间的数据传输安全，提高企业跨区域协同办公能力。同时，通过建设 OA 系统与内部服务器平台，企业能够进一步提升管理效率，降低运营成本。

本方案新增了对\textbf{远程办公人员 VPN 接入}的支持，使外部员工能够通过加密隧道安全访问企业 OA 系统和内部资源，符合后疫情时代混合办公模式的社会需求。

因此，本项目具有较高的社会应用价值和现实意义。

\subsection{技术可行性分析}

本系统主要采用 \textbf{VLAN、VPN（IPSec）、ACL、RIPv2、NAT、DHCP、STP} 等成熟网络技术，具备较高的技术可行性。

\subsubsection{VLAN（虚拟局域网）}

VLAN 技术能够将不同部门划分为独立广播域，实现逻辑隔离，减少广播风暴，提高网络安全性与管理效率。本方案中，总部通过核心交换机 2560-24PS 划分 VLAN，各分公司通过路由器子接口实现 VLAN 隔离。

\subsubsection{VPN（虚拟专用网络）}

VPN 技术用于实现总部与外省分公司之间的安全通信，以及远程办公人员的安全接入。

本方案采用 \textbf{IPSec VPN} 技术，包括：
\begin{itemize}
    \item \textbf{站点到站点 VPN（Site-to-Site VPN）}：总部与三个分公司之间建立 IPSec 隧道，保证跨公网的数据加密传输；
    \item \textbf{远程访问 VPN（Remote Access VPN）}：外部办公 PC 通过 VPN 服务器接入企业内网，安全访问 OA 系统和内部资源。
\end{itemize}

\subsubsection{RIPv2（路由信息协议第二版）}

RIPv2 是一种经典的距离矢量路由协议，支持 VLSM 和 CIDR，适用于中小型网络的动态路由学习。本方案在总部、分公司和 ISP 之间的路由器上部署 RIPv2，实现全网路由自动学习与通告，保证内网可达性。

\subsubsection{ACL（访问控制列表）}

ACL 用于实现精细化的访问控制策略：
\begin{itemize}
    \item 禁止普通部门（如市场部）访问财务部 VLAN；
    \item 仅允许内部员工通过 FTP 协议访问 FTP 服务器；
    \item 限制外部网络对内部 OA 系统的直接访问；
    \item 控制 VPN 流量的加密范围。
\end{itemize}

\subsubsection{NAT（网络地址转换）}

NAT 技术实现内部私有 IP 地址到公网 IP 地址的转换。本方案在 ISP 边界路由器上配置 PAT（端口地址转换），使所有内部终端共享一个公网 IP 地址访问互联网，同时隐藏内部网络拓扑，提高安全性。

\subsubsection{DHCP（动态主机配置协议）}

DHCP 服务器自动为各 VLAN 的终端分配 IP 地址、默认网关和 DNS 服务器地址。本方案将 DHCP 服务与 DNS 服务部署在同一台服务器上（IP: 192.168.80.50），减少人工配置工作量，降低 IP 地址冲突风险。

\subsubsection{STP（生成树协议）}

STP 用于防止交换网络中的环路问题。在总部核心交换机与部门接入交换机之间运行 STP，避免二层环路导致的广播风暴，提高网络稳定性。

\subsection{经济可行性分析}

本项目主要成本包括：

\begin{table}[H]
\centering
\caption{项目预算明细}
\begin{tabular}{|l|r|}
\hline
\textbf{项目}          & \textbf{预算（万元）} \\ \hline
核心路由器（2811）      & 12                   \\ \hline
核心交换机（2560-24PS） & 8                    \\ \hline
接入交换机（2960）      & 18                   \\ \hline
分公司路由器（2811）    & 15                   \\ \hline
分公司交换机（1960-24TT）& 6                    \\ \hline
防火墙设备              & 15                   \\ \hline
服务器设备              & 20                   \\ \hline
VPN 设备与许可          & 10                   \\ \hline
综合布线                & 10                   \\ \hline
施工与维护              & 8                    \\ \hline
\textbf{合计}          & \textbf{约 122}      \\ \hline
\end{tabular}
\end{table}

项目整体预算适中（约 122 万元），符合中小企业信息化建设要求。采用 Packet Tracer 平台进行模拟验证，可在实际采购前充分验证网络设计的合理性与正确性，降低实施风险。

% ==================== 三、需求分析 ====================
\section{需求分析}

\subsection{网络需求}

企业网络需满足以下核心需求：

\begin{itemize}
    \item 支持约 1000 台终端接入，覆盖 7 个部门及 3 个分公司；
    \item 部门之间实现逻辑隔离，财务部等高安全部门禁止普通部门访问；
    \item 所有终端均可通过 NAT 访问互联网；
    \item 总部与三个分公司之间通过 IPSec VPN 实现安全加密互联；
    \item 远程办公人员可通过 VPN 拨入企业内网；
    \item 内部部署 OA 办公系统、邮件系统、WWW 服务和 FTP 服务；
    \item 外部用户可访问企业门户网站和邮件服务；
    \item FTP 服务仅限内部员工使用；
    \item 自动 IP 地址分配（DHCP）；
    \item 网络需具备环路防护（STP）和高扩展性。
\end{itemize}

\subsection{VLAN 规划}

根据企业组织架构和业务需求，将网络划分为 8 个 VLAN，如表 \ref{tab:vlan} 所示。

\begin{table}[H]
\centering
\caption{VLAN 规划表}
\label{tab:vlan}
\begin{tabular}{|c|c|c|c|}
\hline
\textbf{VLAN ID} & \textbf{部门名称} & \textbf{网段}           & \textbf{默认网关}   \\ \hline
VLAN 10          & 市场部 (MARKET)   & 192.168.10.0/24         & 192.168.10.1        \\ \hline
VLAN 20          & 财务部 (FINANCE)  & 192.168.20.0/24         & 192.168.20.1        \\ \hline
VLAN 30          & 人事部 (HR)       & 192.168.30.0/24         & 192.168.30.1        \\ \hline
VLAN 40          & 研发部 (DEV)      & 192.168.40.0/24         & 192.168.40.1        \\ \hline
VLAN 50          & 行政部 (ADMIN)    & 192.168.50.0/24         & 192.168.50.1        \\ \hline
VLAN 60          & 运营部 (OPERATION)& 192.168.60.0/24         & 192.168.60.1        \\ \hline
VLAN 70          & 客服部 (SERVICE)  & 192.168.70.0/24         & 192.168.70.1        \\ \hline
VLAN 80          & 服务器区 (SERVER) & 192.168.80.0/24         & 192.168.80.1        \\ \hline
\end{tabular}
\end{table}

\subsection{服务器规划}

总部服务器区（VLAN 80）部署三台核心服务器，分别承担 Web 服务、DNS/DHCP 服务以及 FTP/Email 服务，如表 \ref{tab:server} 所示。

\begin{table}[H]
\centering
\caption{服务器规划表}
\label{tab:server}
\begin{tabular}{|c|c|c|}
\hline
\textbf{服务器名称} & \textbf{IP 地址}      & \textbf{功能说明}                  \\ \hline
Web Server          & 192.168.80.10         & 企业门户网站（HTTP）               \\ \hline
DNS \& DHCP Server  & 192.168.80.50         & DNS 域名解析 \& DHCP 动态地址分配   \\ \hline
FTP \& Email Server & 192.168.80.20/30      & FTP 文件服务 \& Email 邮件服务      \\ \hline
\end{tabular}
\end{table}

\textbf{设计说明：} 相比于每个服务独立部署一台服务器，本方案将功能相近的服务整合，减少服务器数量，降低硬件和维护成本，同时满足中小企业的实际需求。

\subsection{IP 地址规划}

\begin{table}[H]
\centering
\caption{全局 IP 地址规划}
\label{tab:ip}
\begin{tabular}{|c|c|c|c|}
\hline
\textbf{网络区域}        & \textbf{网段}              & \textbf{子网掩码}       & \textbf{说明}          \\ \hline
总部 VLAN 10--70         & 192.168.10.0/24 -- 70.0/24& 255.255.255.0           & 7 个部门子网           \\ \hline
总部服务器区 (VLAN 80)   & 192.168.80.0/24           & 255.255.255.0           & 服务器集群             \\ \hline
总部与核心路由器互联     & 10.0.0.0/30               & 255.255.255.252         & 点对点链路             \\ \hline
核心路由器与 ISP 互联    & 202.100.1.0/30            & 255.255.255.252         & 公网接口               \\ \hline
分公司 1 内部网段        & 172.16.1.0/24             & 255.255.255.0           & Branch 1               \\ \hline
分公司 2 内部网段        & 172.16.2.0/24             & 255.255.255.0           & Branch 2               \\ \hline
分公司 3 内部网段        & 172.16.3.0/24             & 255.255.255.0           & Branch 3               \\ \hline
分公司与核心路由器互联   & 1.1.1.0/30 等             & 255.255.255.252         & 串行点对点链路         \\ \hline
VPN 客户端地址池         & 172.16.100.0/24           & 255.255.255.0           & 远程 VPN 接入          \\ \hline
\end{tabular}
\end{table}

% ==================== 四、网络结构设计 ====================
\section{网络结构设计}

\subsection{网络拓扑结构}

本方案采用 \textbf{Hub-and-Spoke（中心辐射型）}网络拓扑架构，以总部核心路由器（Core-Router，Cisco 2811）为整个网络的枢纽中心。拓扑结构分为以下层次：

\subsubsection{第一层：ISP 互联层}

ISP 采用 Cisco 2811 路由器，通过串行链路连接核心路由器。ISP 路由器同时连接外部 Web 服务器（模拟互联网服务）和外部测试 PC，用于验证外网访问功能。

\subsubsection{第二层：核心路由层}

核心路由器（Cisco 2811）是整个网络的交通枢纽，负责：
\begin{itemize}
    \item 向上通过串行链路连接 ISP，实现互联网接入；
    \item 向下连接总部顶级路由器（HQ-Router，Cisco 2811）；
    \item 通过三条串行链路分别连接三个分公司路由器，实现分支机构互联；
    \item 运行 RIPv2 动态路由协议，维护全网路由表。
\end{itemize}

\subsubsection{第三层：总部分发层}

总部网络采用两级结构：
\begin{itemize}
    \item \textbf{总部顶级路由器（HQ-Router，Cisco 2811）}：作为总部网络的网关，向上连接核心路由器，向下连接总部核心交换机。
    \item \textbf{总部核心交换机（HQ-Core-SW，Cisco 2560-24PS）}：作为总部网络的核心交换节点，通过 Trunk 链路连接各部门接入交换机和服务器区交换机，支持 VLAN 间路由。
\end{itemize}

\subsubsection{第四层：接入层}

\begin{itemize}
    \item \textbf{部门接入交换机（3 台 Cisco 2960）}：分别连接各部门 PC，承载 VLAN 10--70 的用户流量；
    \item \textbf{服务器接入交换机（1 台 Cisco 2960）}：连接服务器区三台服务器（VLAN 80）；
    \item \textbf{分公司交换机（3 台 Cisco 1960-24TT）}：各分公司内部的二层交换设备，连接分公司终端。
\end{itemize}

\subsubsection{第五层：分支机构}

每个分公司采用\textbf{二级路由结构}：
\begin{itemize}
    \item \textbf{分公司路由器（Cisco 2811）}：通过串行链路与总部核心路由器互联，运行 RIPv2，支持 IPSec VPN 隧道，通过子接口实现 VLAN 划分；
    \item \textbf{分公司交换机（Cisco 1960-24TT）}：连接分公司内部 PC 和终端设备。
\end{itemize}

\subsubsection{拓扑结构图}

整体网络拓扑结构如图 \ref{fig:total-topo} 所示：

\begin{figure}[H]
\centering
\includegraphics[width=\textwidth]{total_topo.png}
\caption{全局网络拓扑图（Cisco Packet Tracer V2）}
\label{fig:total-topo}
\end{figure}

\subsection{VPN 结构设计}

本方案采用双重 VPN 架构，满足不同场景的安全接入需求：

\subsubsection{站点到站点 IPSec VPN}

总部核心路由器与三个分公司路由器之间建立 IPSec VPN 隧道，实现以下功能：
\begin{itemize}
    \item 分公司终端可以安全访问总部的 OA 系统、文件服务器和内部资源；
    \item 所有跨公网传输的数据经过 AES 加密和 SHA 完整性校验；
    \item 使用预共享密钥（Pre-Shared Key）进行身份认证。
\end{itemize}

VPN 对等体关系如下：

\begin{center}
\begin{tabular}{c}
Branch1-Router $\longleftrightarrow$ Core-Router (IPSec Tunnel) \\
Branch2-Router $\longleftrightarrow$ Core-Router (IPSec Tunnel) \\
Branch3-Router $\longleftrightarrow$ Core-Router (IPSec Tunnel) \\
\end{tabular}
\end{center}

\subsubsection{远程访问 VPN}

ISP 侧的外部 PC 通过 VPN 客户端软件，经过公网连接到 VPN 服务器（ISP 路由器上配置），建立加密隧道后安全访问企业内部资源（如 OA 系统）。这是典型的\textbf{Remote Access VPN} 场景，满足远程办公需求。


\subsection{安全设计}

网络安全是本次设计的核心关注点。本方案采用多层次、纵深防御的安全体系：

\begin{enumerate}[label=(\arabic*), leftmargin=2em]
    \item \textbf{VLAN 逻辑隔离}：将不同部门划分到不同的广播域，实现二层隔离，防止部门间的非授权二层访问；
    \item \textbf{ACL 访问控制}：
        \begin{itemize}
            \item 在核心交换机上配置扩展 ACL，禁止普通部门（如市场部、客服部）访问财务部 VLAN 20；
            \item 在服务器区入口配置 ACL，仅允许内部 IP 地址段（192.168.0.0/16）访问 FTP 服务；
            \item 限制外部网络对 OA 系统（TCP 8080）的直接访问，仅允许 VPN 隧道流量通过。
        \end{itemize}
    \item \textbf{NAT 地址转换}：在 ISP 边界路由器上配置 PAT，隐藏内部网络结构，仅暴露一个公网 IP；
    \item \textbf{IPSec VPN 加密}：所有跨公网的站点间和远程接入流量均经过 IPSec 加密，防止数据窃听和篡改；
    \item \textbf{STP 环路防护}：在总部二层交换网络中启用生成树协议，防止广播风暴和 MAC 地址表震荡；
    \item \textbf{路由安全}：RIPv2 支持 MD5 认证（可在后续升级中启用），防止路由欺骗。
\end{enumerate}

% ==================== 五、系统配置与实施 ====================
\section{系统配置与实施}

本章详细介绍各网络设备的关键配置命令。所有配置均基于 Cisco IOS 命令行，并在 Cisco Packet Tracer 平台上完成验证。

\subsection{核心路由器配置（Core-Router）}

设备型号：\textbf{Cisco 2811}，作为全网路由枢纽。

\begin{lstlisting}[caption={核心路由器基础配置}, label={lst:core-router}]
enable
configure terminal

hostname Core-Router

! ---- 连接 ISP 的串行接口 ----
interface Serial0/0/0
 ip address 202.100.1.2 255.255.255.252
 no shutdown

! ---- 连接总部顶级路由器的接口 ----
interface GigabitEthernet0/0
 ip address 10.0.0.1 255.255.255.252
 no shutdown

! ---- 连接分公司 1 的串行接口 ----
interface Serial0/1/0
 ip address 1.1.1.1 255.255.255.252
 no shutdown

! ---- 连接分公司 2 的串行接口 ----
interface Serial0/1/1
 ip address 1.1.2.1 255.255.255.252
 no shutdown

! ---- 连接分公司 3 的串行接口 ----
interface Serial0/2/0
 ip address 1.1.3.1 255.255.255.252
 no shutdown
\end{lstlisting}

\subsection{RIPv2 动态路由配置}

全网采用 \textbf{RIPv2} 动态路由协议，替代旧版方案中的 OSPF。RIPv2 支持 VLSM 和 CIDR，配置更简洁，适合中小型网络。

\begin{lstlisting}[caption={Core-Router RIPv2 路由配置}, label={lst:rip-core}]
router rip
 version 2
 network 202.100.1.0
 network 10.0.0.0
 network 1.1.1.0
 network 1.1.2.0
 network 1.1.3.0
 network 192.168.10.0
 network 192.168.20.0
 network 192.168.30.0
 network 192.168.40.0
 network 192.168.50.0
 network 192.168.60.0
 network 192.168.70.0
 network 192.168.80.0
 no auto-summary
\end{lstlisting}

\begin{lstlisting}[caption={Branch-Router RIPv2 配置（以分公司1为例）}, label={lst:rip-branch}]
! ---- 分公司 1 路由器 (Branch1-Router, 2811) ----
router rip
 version 2
 network 1.1.1.0
 network 172.16.1.0
 no auto-summary
\end{lstlisting}

\begin{lstlisting}[caption={总部顶级路由器 RIPv2 配置}, label={lst:rip-hq}]
! ---- 总部顶级路由器 (HQ-Router, 2811) ----
router rip
 version 2
 network 10.0.0.0
 network 192.168.10.0
 network 192.168.20.0
 network 192.168.30.0
 network 192.168.40.0
 network 192.168.50.0
 network 192.168.60.0
 network 192.168.70.0
 network 192.168.80.0
 no auto-summary
\end{lstlisting}

\subsection{总部核心交换机配置（HQ-Core-SW）}

设备型号：\textbf{Cisco 2560-24PS}。

\begin{lstlisting}[caption={总部核心交换机 VLAN 创建}, label={lst:core-sw-vlan}]
enable
configure terminal

hostname HQ-Core-SW

! ---- 创建 VLAN ----
vlan 10
 name MARKET
vlan 20
 name FINANCE
vlan 30
 name HR
vlan 40
 name DEV
vlan 50
 name ADMIN
vlan 60
 name OPERATION
vlan 70
 name SERVICE
vlan 80
 name SERVER
\end{lstlisting}

\begin{lstlisting}[caption={总部核心交换机 SVI 接口配置}, label={lst:core-sw-svi}]
! ---- 配置各 VLAN 的三层 SVI 接口（VLAN 间路由） ----
interface vlan 10
 ip address 192.168.10.1 255.255.255.0
 no shutdown
!
interface vlan 20
 ip address 192.168.20.1 255.255.255.0
 no shutdown
!
interface vlan 30
 ip address 192.168.30.1 255.255.255.0
 no shutdown
!
interface vlan 40
 ip address 192.168.40.1 255.255.255.0
 no shutdown
!
interface vlan 50
 ip address 192.168.50.1 255.255.255.0
 no shutdown
!
interface vlan 60
 ip address 192.168.60.1 255.255.255.0
 no shutdown
!
interface vlan 70
 ip address 192.168.70.1 255.255.255.0
 no shutdown
!
interface vlan 80
 ip address 192.168.80.1 255.255.255.0
 no shutdown
\end{lstlisting}

\begin{lstlisting}[caption={Trunk 端口配置}, label={lst:core-sw-trunk}]
! ---- 配置与部门接入交换机和服务器交换机之间的 Trunk 链路 ----
interface range FastEthernet0/1-4
 switchport mode trunk

! ---- 配置与顶级路由器连接的上行端口 ----
interface GigabitEthernet0/1
 no switchport
 ip address 10.0.0.2 255.255.255.252
 no shutdown
\end{lstlisting}

\subsection{DHCP 配置}

DHCP 服务由服务器区的 DNS/DHCP 服务器（192.168.80.50）提供，同时在核心交换机上配置 DHCP 中继（如需）。以下为服务器端 DHCP 配置示意和各 VLAN 地址池配置：

\begin{lstlisting}[caption={核心交换机 DHCP 配置}, label={lst:dhcp}]
! ---- 排除静态 IP 地址范围 ----
ip dhcp excluded-address 192.168.10.1 192.168.10.20
ip dhcp excluded-address 192.168.20.1 192.168.20.20
ip dhcp excluded-address 192.168.30.1 192.168.30.20

! ---- 各 VLAN DHCP 地址池 ----
ip dhcp pool MARKET
 network 192.168.10.0 255.255.255.0
 default-router 192.168.10.1
 dns-server 192.168.80.50

ip dhcp pool FINANCE
 network 192.168.20.0 255.255.255.0
 default-router 192.168.20.1
 dns-server 192.168.80.50

ip dhcp pool HR
 network 192.168.30.0 255.255.255.0
 default-router 192.168.30.1
 dns-server 192.168.80.50

! ---- 其余 VLAN 地址池配置同理 ----
\end{lstlisting}

\subsection{边界路由器与 NAT 配置}

设备型号：\textbf{Cisco 2811（ISP-Router 和 Core-Router）}。

\begin{lstlisting}[caption={NAT 配置（Core-Router）}, label={lst:nat}]
! ---- 定义内外网接口 ----
interface Serial0/0/0
 ip nat outside

interface GigabitEthernet0/0
 ip nat inside

! ---- 定义 NAT 转换 ACL ----
access-list 1 permit 192.168.0.0 0.0.255.255
access-list 1 permit 172.16.0.0 0.0.255.255

! ---- PAT 转换（端口地址转换） ----
ip nat inside source list 1 interface Serial0/0/0 overload

! ---- 默认路由指向 ISP ----
ip route 0.0.0.0 0.0.0.0 202.100.1.1
\end{lstlisting}

\subsection{ACL 访问控制配置}

\begin{lstlisting}[caption={ACL 配置 — 部门间访问控制}, label={lst:acl-dept}]
! ---- 禁止市场部（VLAN 10）访问财务部（VLAN 20） ----
access-list 100 deny ip 192.168.10.0 0.0.0.255 192.168.20.0 0.0.0.255
access-list 100 deny ip 192.168.70.0 0.0.0.255 192.168.20.0 0.0.0.255
access-list 100 permit ip any any
!
interface vlan 10
 ip access-group 100 in
interface vlan 70
 ip access-group 100 in
\end{lstlisting}

\begin{lstlisting}[caption={ACL 配置 — FTP 访问控制}, label={lst:acl-ftp}]
! ---- 仅允许内部 IP 访问 FTP 服务器（192.168.80.20） ----
access-list 110 permit tcp 192.168.0.0 0.0.255.255 host 192.168.80.20 eq 21
access-list 110 permit tcp 172.16.0.0 0.0.255.255 host 192.168.80.20 eq 21
access-list 110 deny tcp any host 192.168.80.20 eq 21
access-list 110 permit ip any any
!
interface vlan 80
 ip access-group 110 out
\end{lstlisting}

\subsection{IPSec VPN 配置}

\subsubsection{站点到站点 VPN 配置（Core-Router 侧）}

\begin{lstlisting}[caption={IPSec VPN Phase 1 — IKE/ISAKMP 策略}, label={lst:vpn-ike}]
! ---- IKE Phase 1 策略 ----
crypto isakmp policy 10
 encryption aes 256
 hash sha
 authentication pre-share
 group 2
 lifetime 86400

! ---- 预共享密钥（分别对应三个分公司） ----
crypto isakmp key Branch1Key address 1.1.1.2
crypto isakmp key Branch2Key address 1.1.2.2
crypto isakmp key Branch3Key address 1.1.3.2
\end{lstlisting}

\begin{lstlisting}[caption={IPSec VPN Phase 2 — 转换集与加密映射}, label={lst:vpn-ipsec}]
! ---- IPSec 转换集 ----
crypto ipsec transform-set VPN-SET esp-aes 256 esp-sha-hmac

! ---- 定义 VPN 加密流量 ----
access-list 120 permit ip 192.168.0.0 0.0.255.255 172.16.1.0 0.0.0.255
access-list 121 permit ip 192.168.0.0 0.0.255.255 172.16.2.0 0.0.0.255
access-list 122 permit ip 192.168.0.0 0.0.255.255 172.16.3.0 0.0.0.255

! ---- 加密映射（分别对应三个分公司） ----
crypto map VPN-MAP 10 ipsec-isakmp
 set peer 1.1.1.2
 set transform-set VPN-SET
 match address 120

crypto map VPN-MAP 20 ipsec-isakmp
 set peer 1.1.2.2
 set transform-set VPN-SET
 match address 121

crypto map VPN-MAP 30 ipsec-isakmp
 set peer 1.1.3.2
 set transform-set VPN-SET
 match address 122

! ---- 将加密映射应用到接口 ----
interface Serial0/1/0
 crypto map VPN-MAP
interface Serial0/1/1
 crypto map VPN-MAP
interface Serial0/2/0
 crypto map VPN-MAP
\end{lstlisting}

\subsubsection{远程访问 VPN 配置（ISP-Router/VPN-Server）}

\begin{lstlisting}[caption={远程访问 VPN 配置}, label={lst:vpn-remote}]
! ---- VPN 服务器（ISP-Router）上的远程接入配置 ----
username remote_user password 0 SecurePass123

! ---- 为远程 VPN 客户端分配地址池 ----
ip local pool VPN-POOL 172.16.100.10 172.16.100.50

! ---- 远程接入 VPN 策略 ----
crypto isakmp client configuration group VPN-GROUP
 key RemoteVPNKey
 pool VPN-POOL
 acl 130

access-list 130 permit ip 192.168.0.0 0.0.255.255 any
\end{lstlisting}

\subsection{服务器配置}

\subsubsection{WEB 服务器（192.168.80.10）}

\begin{itemize}
    \item \textbf{HTTP 服务}：开启，端口 80
    \item \textbf{首页内容}：企业门户网站欢迎页
\begin{lstlisting}[language=HTML]
<html>
<head><title>企业门户网站</title></head>
<body>
  <h1>欢迎访问企业官方网站</h1>
  <p>提供产品展示、企业动态、在线客服等服务</p>
</body>
</html>
\end{lstlisting}
\end{itemize}

\subsubsection{FTP \& Email 服务器（192.168.80.20/30）}

\begin{itemize}
    \item \textbf{FTP 服务}：开启，端口 21，支持匿名登录（Anonymous Login: Enable），通过 ACL 限制仅内部员工访问
    \item \textbf{Email 服务}：开启，支持 SMTP（端口 25）和 POP3（端口 110），用于企业内部和外部邮件通信
\end{itemize}

\subsubsection{DNS \& DHCP 服务器（192.168.80.50）}

\begin{itemize}
    \item \textbf{DNS 服务}：提供域名解析，将 \texttt{www.company.com} 解析到 192.168.80.10，\texttt{mail.company.com} 解析到 192.168.80.30
    \item \textbf{DHCP 服务}：为各 VLAN 终端自动分配 IP、网关、DNS
\end{itemize}

\subsection{分公司路由器配置（Branch-Router）}

设备型号：\textbf{Cisco 2811}（每个分公司各一台）。

\begin{lstlisting}[caption={分公司路由器配置（以Branch1为例）}, label={lst:branch}]
enable
configure terminal

hostname Branch1-Router

! ---- 连接核心路由器的串行接口 ----
interface Serial0/0/0
 ip address 1.1.1.2 255.255.255.252
 no shutdown

! ---- 连接分公司交换机的接口 ----
interface GigabitEthernet0/0
 ip address 172.16.1.1 255.255.255.0
 no shutdown

! ---- RIPv2 路由 ----
router rip
 version 2
 network 1.1.1.0
 network 172.16.1.0
 no auto-summary

! ---- IPSec VPN（对应核心路由器的配置） ----
crypto isakmp policy 10
 encryption aes 256
 hash sha
 authentication pre-share
 group 2
crypto isakmp key Branch1Key address 1.1.1.1
crypto ipsec transform-set VPN-SET esp-aes 256 esp-sha-hmac
access-list 120 permit ip 172.16.1.0 0.0.0.255 192.168.0.0 0.0.255.255
crypto map VPN-MAP 10 ipsec-isakmp
 set peer 1.1.1.1
 set transform-set VPN-SET
 match address 120
interface Serial0/0/0
 crypto map VPN-MAP
\end{lstlisting}

\subsection{分公司交换机配置（Branch-SW）}

设备型号：\textbf{Cisco 1960-24TT}（每个分公司各一台）。

\begin{lstlisting}[caption={分公司交换机基础配置}, label={lst:branch-sw}]
enable
configure terminal

hostname Branch1-SW

! ---- 创建本地 VLAN ----
vlan 100
 name Branch1-LAN

! ---- 配置管理 IP ----
interface vlan 100
 ip address 172.16.1.2 255.255.255.0
 no shutdown

! ---- 用户端口 ----
interface range FastEthernet0/1-24
 switchport access vlan 100
 no shutdown

! ---- 上行端口（连接分公司路由器） ----
interface GigabitEthernet0/1
 switchport mode trunk
 no shutdown
\end{lstlisting}

% ==================== 六、设备选型与工程预算 ====================
\section{设备选型与工程预算}

\subsection{核心网络设备清单}

本方案所有设备均可在 Cisco Packet Tracer 8.x 平台上仿真运行。

\begin{table}[H]
\centering
\caption{核心网络设备清单}
\label{tab:device-list}
\begin{tabular}{|c|c|c|c|}
\hline
\textbf{设备名称}       & \textbf{型号}           & \textbf{数量} & \textbf{部署位置}   \\ \hline
核心路由器              & Cisco 2811              & 1             & 总部中心机房         \\ \hline
总部顶级路由器          & Cisco 2811              & 1             & 总部中心机房         \\ \hline
ISP 路由器/VPN 服务器   & Cisco 2811              & 1             & ISP 机房             \\ \hline
分公司路由器            & Cisco 2811              & 3             & 各分公司机房         \\ \hline
总部核心交换机          & Cisco 2560-24PS         & 1             & 总部中心机房         \\ \hline
部门接入交换机          & Cisco 2960              & 3             & 总部各楼层弱电间     \\ \hline
服务器接入交换机        & Cisco 2960              & 1             & 总部中心机房         \\ \hline
分公司交换机            & Cisco 1960-24TT         & 3             & 各分公司机房         \\ \hline
Web 服务器              & Server-PT（模拟）       & 1             & 总部服务器区         \\ \hline
DNS/DHCP 服务器         & Server-PT（模拟）       & 1             & 总部服务器区         \\ \hline
FTP/Email 服务器        & Server-PT（模拟）       & 1             & 总部服务器区         \\ \hline
\end{tabular}
\end{table}

\subsection{预算测算依据}

\subsubsection{接入端口测算}

1000 台 PC，每台接入交换机（Cisco 2960）提供 48 个端口，则：
\[
\text{所需交换机数量} = \lceil 1000 / 48 \rceil = \lceil 20.83 \rceil = 21 \text{ 台}
\]

考虑 3 个部门接入交换机 + 1 台服务器交换机共 4 台覆盖总部核心接入，其余终端通过分公司交换机（3 台 1960-24TT，每台 24 口）承载。

\subsubsection{核心带宽测算}

假设单用户平均带宽需求为 20 Mbps，并发率为 30\%，则核心交换机上行带宽需求为：
\[
\text{上行带宽} = 1000 \times 20\text{ Mbps} \times 30\% = 6000\text{ Mbps} = 6\text{ Gbps}
\]

核心交换机（2560-24PS）提供 24 个千兆端口，支持线速转发，满足带宽需求。

\subsubsection{项目总预算}

\begin{table}[H]
\centering
\caption{项目总预算}
\begin{tabular}{|l|r|r|}
\hline
\textbf{设备/项目}       & \textbf{单价（万元）} & \textbf{总价（万元）} \\ \hline
核心路由器 2811 × 1      & 5                     & 5                     \\ \hline
总部分发路由器 2811 × 1  & 5                     & 5                     \\ \hline
ISP 路由器 2811 × 1      & 5                     & 5                     \\ \hline
分公司路由器 2811 × 3    & 5                     & 15                    \\ \hline
核心交换机 2560-24PS × 1 & 8                     & 8                     \\ \hline
接入交换机 2960 × 4      & 4.5                   & 18                    \\ \hline
分公司交换机 1960-24TT × 3 & 2                   & 6                     \\ \hline
服务器 × 3               & 6.67                  & 20                    \\ \hline
综合布线                  & --                    & 10                    \\ \hline
VPN 设备与许可            & --                    & 10                    \\ \hline
施工与维护                & --                    & 8                     \\ \hline
\textbf{总计}            &                       & \textbf{110}          \\ \hline
\end{tabular}
\end{table}

项目总预算约为 110 万元，符合中小企业信息化建设预算范围。

% ==================== 七、构建和测试 ====================
\section{构建和测试}

本章展示基于 Cisco Packet Tracer 平台搭建的完整网络拓扑及各项功能测试结果。

\subsection{全局拓扑图}

完整网络拓扑已在 Cisco Packet Tracer 中成功搭建并运行。如图 \ref{fig:total-topo} 所示，所有设备连接正确，路由协议正常工作。

\subsection{内网 VLAN 通信测试}

\subsubsection{同一 VLAN 内通信}

测试同一 VLAN（VLAN 10 市场部）内两台 PC 之间的通信。如图 \ref{fig:inside-vlan} 所示，同 VLAN 内的终端可以正常 ping 通，二层转发正常。

\begin{figure}[H]
\centering
\includegraphics[width=0.85\textwidth]{inside_vlan_free.png}
\caption{同 VLAN 内通信测试——通信正常}
\label{fig:inside-vlan}
\end{figure}

\subsubsection{同公司不同 VLAN 间通信}

测试同一公司内部不同 VLAN 之间（如 VLAN 10 市场部与 VLAN 80 服务器区）的通信。通过三层路由（SVI 接口），不同 VLAN 间可以正常通信，如图 \ref{fig:between-vlan-inside} 所示。

\begin{figure}[H]
\centering
\includegraphics[width=0.85\textwidth]{between_vlan_inside_comp_free.png}
\caption{同公司 VLAN 间通信测试——通信正常}
\label{fig:between-vlan-inside}
\end{figure}

\subsubsection{不同公司间通信}

测试总部与分公司之间的通信。通过 RIPv2 动态路由和 IPSec VPN 隧道，总部与分公司之间通信正常（图 \ref{fig:between-vlan-comp-free}）。

测试总部不同部门与分公司之间的通信（图 \ref{fig:between-vlan-comp-free}），验证 VLAN 间路由和跨站点路由的正确性。

\begin{figure}[H]
\centering
\includegraphics[width=0.85\textwidth]{between_vlan_between_comp_free.png}
\caption{总部与分公司 VLAN 间通信——通信正常}
\label{fig:between-vlan-comp-free}
\end{figure}

\subsection{ACL 访问控制测试}

验证 VLAN 10（市场部）访问 VLAN 20（财务部）被 ACL 阻止。如图 \ref{fig:acl-control} 所示，ACL 规则生效，市场部 PC 无法 ping 通财务部 PC，访问控制策略正常执行。

\begin{figure}[H]
\centering
\includegraphics[width=0.85\textwidth]{between_vlan_between_comp_control.png}
\caption{ACL 访问控制测试——市场部无法访问财务部}
\label{fig:acl-control}
\end{figure}

\subsection{DHCP 自动地址分配测试}

如图 \ref{fig:dhcp} 所示，终端通过 DHCP 成功获取 IP 地址、默认网关和 DNS 服务器地址，验证 DHCP 服务配置正确。

\begin{figure}[H]
\centering
\includegraphics[width=0.85\textwidth]{dhcp_allocate_ip.png}
\caption{DHCP 动态 IP 地址分配测试}
\label{fig:dhcp}
\end{figure}

\subsection{互联网访问测试}

通过 NAT/PAT 技术，内网终端可以正常访问互联网。如图 \ref{fig:internet} 所示，总部 PC 成功 ping 通外部服务器，并可通过 HTTP 访问外部 Web 服务，验证 NAT 转换功能正常。

\begin{figure}[H]
\centering
\includegraphics[width=0.85\textwidth]{internet_funciton.png}
\caption{内网终端访问互联网测试}
\label{fig:internet}
\end{figure}

\subsection{FTP 服务测试}

内部 PC 通过 FTP 协议成功连接并登录 FTP 服务器（192.168.80.20），可以进行文件上传和下载操作，如图 \ref{fig:ftp} 所示。同时验证外部 PC 无法访问 FTP 服务（ACL 限制），满足安全需求。

\begin{figure}[H]
\centering
\includegraphics[width=0.85\textwidth]{ftp_function.png}
\caption{FTP 服务功能测试}
\label{fig:ftp}
\end{figure}

\subsection{Email 邮件服务测试}

如图 \ref{fig:email} 所示，内部用户之间以及内外部用户之间的邮件收发功能正常，验证 Email 服务器的 SMTP 和 POP3 服务正常工作。

\begin{figure}[H]
\centering
\includegraphics[width=0.85\textwidth]{email_function.png}
\caption{Email 邮件服务功能测试}
\label{fig:email}
\end{figure}

\subsection{VPN 功能测试}

\subsubsection{内网 Ping VPN 服务器}

总部内网 PC 可以成功 ping 通 VPN 服务器（ISP-Router 的内部接口），验证到达 VPN 服务器的路由可达，如图 \ref{fig:ping-vpn-server} 所示。

\begin{figure}[H]
\centering
\includegraphics[width=0.85\textwidth]{inner_ping_vpn_server_success.png}
\caption{内网 Ping VPN 服务器测试——通信正常}
\label{fig:ping-vpn-server}
\end{figure}

\subsubsection{VPN 连接建立}

远程外部 PC 通过 VPN 客户端成功与 VPN 服务器建立 IPSec 加密隧道连接，如图 \ref{fig:vpn-establish} 所示。隧道建立后，远程 PC 获得虚拟内网 IP 地址（172.16.100.0/24 地址池），可以安全访问企业内网资源。

\begin{figure}[H]
\centering
\includegraphics[width=0.85\textwidth]{vpn_connect_establish.png}
\caption{VPN 隧道建立测试——连接成功}
\label{fig:vpn-establish}
\end{figure}

\subsubsection{VPN 连接后通信测试}

VPN 隧道建立后，外部 PC 通过加密隧道成功 ping 通企业内网服务器，如图 \ref{fig:vpn-ping} 所示，验证 VPN 隧道的端到端可达性。

\begin{figure}[H]
\centering
\includegraphics[width=0.85\textwidth]{vpn_connect_ping_succcess.png}
\caption{VPN 隧道通信测试——Ping 成功}
\label{fig:vpn-ping}
\end{figure}

\subsection{测试结果汇总}

\begin{table}[H]
\centering
\caption{功能测试结果汇总}
\label{tab:test-summary}
\begin{tabular}{|c|c|c|c|}
\hline
\textbf{测试项目}           & \textbf{测试方法}    & \textbf{预期结果}   & \textbf{实际结果}   \\ \hline
同 VLAN 内通信              & Ping                 & 通信正常            & $\checkmark$ 通过    \\ \hline
不同 VLAN 间通信            & Ping                 & 通信正常            & $\checkmark$ 通过    \\ \hline
ACL 部门间访问控制          & Ping（被拒绝）       & 无法通信            & $\checkmark$ 通过    \\ \hline
DHCP 地址分配               & Ipconfig             & 正确获取 IP         & $\checkmark$ 通过    \\ \hline
NAT 互联网访问              & Ping + HTTP          & 正常访问外网        & $\checkmark$ 通过    \\ \hline
FTP 服务（内部）            & FTP 登录             & 正常访问            & $\checkmark$ 通过    \\ \hline
FTP 服务（外部）            & FTP 登录（被拒绝）   & 无法访问            & $\checkmark$ 通过    \\ \hline
Email 邮件服务              & 发送/接收邮件        & 正常收发            & $\checkmark$ 通过    \\ \hline
内网 Ping VPN 服务器        & Ping                 & 通信正常            & $\checkmark$ 通过    \\ \hline
VPN 隧道建立                & VPN 拨号连接         & 连接成功            & $\checkmark$ 通过    \\ \hline
VPN 隧道通信                & Ping（隧道内）       & 通信正常            & $\checkmark$ 通过    \\ \hline
RIPv2 路由学习              & show ip route        & 路由表完整          & $\checkmark$ 通过    \\ \hline
分公司与总部互联            & Ping + HTTP          & 双向通信正常        & $\checkmark$ 通过    \\ \hline
\end{tabular}
\end{table}

所有测试项目均通过，验证了本方案设计的正确性和可行性。

% ==================== 八、总结与展望 ====================
\section{总结与展望}

\subsection{项目总结}

本方案采用 Cisco Packet Tracer 平台，基于 \textbf{Hub-and-Spoke} 网络拓扑架构，完成了中小企业网络规划与设计。方案综合运用了以下关键网络技术：

\begin{enumerate}[label=(\arabic*), leftmargin=2em]
    \item \textbf{VLAN 逻辑隔离}：将 7 个部门和服务器区划分为 8 个独立广播域，实现二层隔离和三层互通；
    \item \textbf{RIPv2 动态路由}：替代旧版方案中的 OSPF，使用更简洁的距离矢量路由协议实现全网路由自动学习；
    \item \textbf{ACL 访问控制}：实现部门间精细化访问控制和 FTP 服务的内部限用策略；
    \item \textbf{NAT 地址转换}：通过 PAT 实现内网终端共享公网 IP 访问互联网；
    \item \textbf{IPSec VPN}：实现总部与三个分公司的站点到站点加密通信，以及远程办公人员的安全接入；
    \item \textbf{DHCP 自动编址}：实现终端 IP 地址的自动分配和管理；
    \item \textbf{STP 环路防护}：确保二层网络的稳定运行。
\end{enumerate}

\subsection{方案特色}

与旧版方案相比，本 V2 方案具有以下改进和特色：

\begin{itemize}
    \item \textbf{更清晰的层次结构}：采用 Hub-and-Spoke 架构，核心路由器作为全网枢纽，层次分明，便于管理和故障排查；
    \item \textbf{真实的设备选型}：使用 Cisco Packet Tracer 中具体可仿真的设备型号（2811 路由器、2560-24PS 交换机、2960 交换机、1960-24TT 交换机）；
    \item \textbf{双重 VPN 架构}：同时支持站点到站点 VPN 和远程访问 VPN，满足更多业务场景；
    \item \textbf{服务整合}：将 DNS 与 DHCP 整合、FTP 与 Email 整合，减少服务器数量，降低运维成本；
    \item \textbf{更全面的测试覆盖}：包含 13 项功能测试，覆盖 VLAN、ACL、DHCP、NAT、FTP、Email、VPN 等全部功能点；
    \item \textbf{实用优先的设计理念}：RIPv2 相比 OSPF 在小网络中配置更简单，更适合中小企业的实际运维能力。
\end{itemize}

\subsection{未来展望}

本方案虽然已满足当前需求，但仍可在以下方面进行改进和扩展：

\begin{enumerate}[label=(\arabic*), leftmargin=2em]
    \item \textbf{高可用性}：引入 HSRP/VRRP 网关冗余协议，部署双核心路由器和双核心交换机，消除单点故障；
    \item \textbf{安全增强}：部署独立的防火墙设备（如 Cisco ASA），在 DMZ 区部署对外服务器，实现更严格的安全隔离；
    \item \textbf{QoS 质量保证}：引入 QoS 策略，为 VoIP、视频会议等实时业务提供带宽保障；
    \item \textbf{无线网络覆盖}：增加无线 AP 和无线控制器（WLC），支持移动办公终端接入；
    \item \textbf{IPv6 支持}：引入 IPv6 双栈部署，为未来网络演进做好准备；
    \item \textbf{网络监控}：部署 SNMP 网络管理系统（如 Cisco Prime），实现全网设备的集中监控和故障告警。
\end{enumerate}

\vspace{1cm}

\begin{flushright}
\begin{minipage}{6cm}
\centering
\rule{6cm}{0.4pt}\\[0.3cm]
\textbf{参考文献}

{[1]} Cisco Systems, Inc. \textit{Cisco IOS IP Routing: RIP Configuration Guide}. Cisco Documentation, 2024.

{[2]} William Stallings. \textit{Data and Computer Communications}. 10th Edition, Pearson, 2019.

{[3]} James F. Kurose, Keith W. Ross. \textit{Computer Networking: A Top-Down Approach}. 8th Edition, Pearson, 2020.

{[4]} Cisco Networking Academy. \textit{CCNA Routing and Switching: Connecting Networks}. Cisco Press, 2023.

{[5]} 谢希仁. \textit{计算机网络}. 第8版, 电子工业出版社, 2021.
\end{minipage}
\end{flushright}

\end{document}
